\chapter{证据质量与可审计边界}

\section{行情与股本}

16只初始候选的2026年7月22日收盘价均经AStock/Sina与腾讯行情双源匹配。六只深度标的的总股本均与公司公告或季度报告对齐。恩捷股份使用限制性股票注销后的9.8132亿股，江波龙使用限制性股票归属后的4.2306亿股。

行情质量足以确认现价与市值，但不能识别卖方身份、融资渠道或下跌原因。近十日回撤只用于刻画估值压缩，不自动解释为基本面恶化。

\section{财务与预告}

六只深度标的的2025年报与2026Q1收入、归母、扣非与经营现金流均与官方报告核对。H1预告直接来自公司公告，但均未经审计；除江波龙披露H1收入区间外，其余公司没有在预告中披露收入。锂盐与隔膜公司也未披露H1吨位、ASP、单位成本或单平利润。

\begin{exhibitbox}[Exhibit 13：官方财务锚]
\small
\begin{tabularx}{\exhibitboxwidth}{L{2.4cm} R{1.7cm} R{1.7cm} R{1.7cm} R{1.7cm} X}
\toprule
标的 & 2025收入 & Q1归母 & Q1 OCF & H1归母预告 & 关键质量结论 \\
\midrule
中国船舶 & 1,519.78亿 & 48.32亿 & +81.52亿 & 92--110亿 & 并表重述但现金强 \\
江波龙 & 227.66亿 & 38.62亿 & -28.75亿 & 92--110亿 & 利润/现金严重背离 \\
恩捷股份 & 136.33亿 & 2.60亿 & +2.07亿 & 7.36--9.00亿 & 扭亏已确认，单位经济未披露 \\
盛新锂能 & 50.64亿 & 4.64亿 & -6.67亿 & 10--12亿 & 周期反转，现金待验证 \\
天华新能 & 75.49亿 & 9.69亿 & -1.83亿 & 22--24亿 & 利润强，现金与资源自供待验证 \\
雅化集团 & 85.43亿 & 3.39亿 & -4.31亿 & 11--13亿 & 公司证据完整，现金仍是阻断 \\
\bottomrule
\end{tabularx}
\sourcenote{公司2025年报、2026Q1报告、2026H1预告。H1预告未经审计。}
\end{exhibitbox}

\section{券商证据}

本次归档14份原始PDF：中国船舶5份、江波龙2份、恩捷股份1份、盛新锂能4份、天华新能1份、雅化集团1份。只有雅化和恩捷的报告同时披露目标、预测与估值方法。其他报告的盈利预测仍可比较，但目标价权重为零。

雅化目标42元对应2026E EPS 2.63元乘16倍PE；恩捷目标103元对应2027E EPS 5.16元乘20倍PE。两份报告的目标期限均未披露，因而不支持“卖方预计年末翻倍”的表述。

\section{事实、预测与推断}

\begin{exhibitbox}[Exhibit 14：证据金字塔]
\centering
\begin{tikzpicture}[font=\small]
\fill[navy!85] (-1.7,3.0)--(1.7,3.0)--(0,4.1)--cycle;
\node[text=white,align=center] at (0,3.45) {公司公告\\已披露事实};
\fill[accentblue!75] (-3.0,1.9)--(3.0,1.9)--(1.7,3.0)--(-1.7,3.0)--cycle;
\node[text=white,align=center] at (0,2.45) {原始券商PDF\\预测、方法、目标};
\fill[steelgrey!65] (-4.2,0.8)--(4.2,0.8)--(3.0,1.9)--(-3.0,1.9)--cycle;
\node[text=white,align=center] at (0,1.35) {实时行情与结构化数据\\现价、市值、历史交易};
\fill[mutedgrey!35] (-5.2,-0.4)--(5.2,-0.4)--(4.2,0.8)--(-4.2,0.8)--cycle;
\node[align=center] at (0,0.2) {AStock推断\\概率、验证阈值、年末适用性};
\end{tikzpicture}
\end{exhibitbox}

事实可以直接陈述；预测必须注明来源与年份；推断必须说明假设和失效条件。下游需求、客户关系、平台认证、名义产能与股权绑定都不能独立证明订单、收入、利用率或利润。

\section{仍然没有解决的缺口}

截至截止日，未取得六只标的统一口径的当日官方北向持股、两融与席位数据，因此不以聚合数据替代。没有取得中国船舶、江波龙、盛新锂能与天华新能的显式当前外部目标；没有取得雅化与恩捷目标价的明确期限；恩捷客户涨价、实际单平利润与有效利用率仍未披露。

这些缺口不是报告“没做完”，而是公开证据边界。正确处理方式是零权、降级与持续验证，而不是填入未经证实的数字。

\clearpage
